\appendix
\section{ALU代码 (alu.v) —— always实现}

参考指导书Listing 7的always实现方式，32位ALU模块支持6种运算（含SLT）：

\begin{lstlisting}[language=Verilog]
module alu (
    input  [31:0] A,
    input  [31:0] B,
    input  [2:0]  F,
    output reg [31:0] Result
);

    always @(*) begin
        case (F)
            3'b000: Result = A + B;
            3'b001: Result = A - B;
            3'b010: Result = A & B;
            3'b011: Result = A | B;
            3'b100: Result = ~A;
            3'b101: Result = ($signed(A) < $signed(B)) ? 32'd1 : 32'd0;
            default: Result = 32'd0;
        endcase
    end

endmodule
\end{lstlisting}

\section{ALU代码 (alu\_assign.v) —— assign实现}

参考指导书Listing 6的assign实现方式（不含SLT，仅5种基本运算）：

\begin{lstlisting}[language=Verilog]
module alu_assign (
    input  wire [7:0]  num1,
    input  wire [2:0]  op,
    output wire [31:0] result
);

    wire [31:0] num2;
    wire [31:0] Sign_extend;

    assign num2 = 32'h00000001;
    assign Sign_extend = {24'b0, num1[7:0]};

    assign result = (op == 3'b000) ? Sign_extend + num2 :
                    (op == 3'b001) ? Sign_extend - num2 :
                    (op == 3'b010) ? Sign_extend & num2 :
                    (op == 3'b011) ? Sign_extend | num2 :
                    (op == 3'b100) ? ~Sign_extend :
                    32'h00000000;

endmodule
\end{lstlisting}

\section{4级流水线8bit全加器代码 (adder\_8bits\_4steps.v)}

参考指导书Listing 5的4级流水线8位加法器，增加stall和flush控制：

\begin{lstlisting}[language=Verilog]
module adder_8bits_4steps (
    input         clk,
    input         rst_n,
    input  [7:0]  cin_a,
    input  [7:0]  cin_b,
    input         c_in,
    input         stall,
    input         flush,
    output        c_out,
    output [7:0]  sum_out
);

    reg        c_out;
    reg        c_out_t1, c_out_t2, c_out_t3;
    reg [7:0]  sum_out;
    reg [1:0]  sum_out_t1;
    reg [3:0]  sum_out_t2;
    reg [5:0]  sum_out_t3;

    // Stage 1: bits[1:0]
    always @(posedge clk or negedge rst_n) begin
        if (!rst_n) begin
            {c_out_t1, sum_out_t1} <= 3'd0;
        end else if (flush) begin
            {c_out_t1, sum_out_t1} <= 3'd0;
        end else if (!stall) begin
            {c_out_t1, sum_out_t1} <= {1'b0, cin_a[1:0]} + {1'b0, cin_b[1:0]} + c_in;
        end
    end

    // Stage 2: bits[3:2]
    always @(posedge clk or negedge rst_n) begin
        if (!rst_n) begin
            {c_out_t2, sum_out_t2} <= 5'd0;
        end else if (flush) begin
            {c_out_t2, sum_out_t2} <= 5'd0;
        end else if (!stall) begin
            {c_out_t2, sum_out_t2} <= {{1'b0, cin_a[3:2]} + {1'b0, cin_b[3:2]} + c_out_t1, sum_out_t1};
        end
    end

    // Stage 3: bits[5:4]
    always @(posedge clk or negedge rst_n) begin
        if (!rst_n) begin
            {c_out_t3, sum_out_t3} <= 7'd0;
        end else if (flush) begin
            {c_out_t3, sum_out_t3} <= 7'd0;
        end else if (!stall) begin
            {c_out_t3, sum_out_t3} <= {{1'b0, cin_a[5:4]} + {1'b0, cin_b[5:4]} + c_out_t2, sum_out_t2};
        end
    end

    // Stage 4: bits[7:6], final output
    always @(posedge clk or negedge rst_n) begin
        if (!rst_n) begin
            {c_out, sum_out} <= 9'd0;
        end else if (!stall) begin
            {c_out, sum_out} <= {{1'b0, cin_a[7:6]} + {1'b0, cin_b[7:6]} + c_out_t3, sum_out_t3};
        end
    end

endmodule
\end{lstlisting}

\section{流水线全加器测试激励代码 (tb\_pipeline\_adder.v)}

\begin{lstlisting}[language=Verilog]
module tb_pipeline_adder;
    reg        clk;
    reg        rst_n;
    reg  [7:0] cin_a;
    reg  [7:0] cin_b;
    reg        c_in;
    reg        stall;
    reg        flush;
    wire       c_out;
    wire [7:0] sum_out;

    integer cycle_count;

    adder_8bits_4steps uut (
        .clk(clk),
        .rst_n(rst_n),
        .cin_a(cin_a),
        .cin_b(cin_b),
        .c_in(c_in),
        .stall(stall),
        .flush(flush),
        .c_out(c_out),
        .sum_out(sum_out)
    );

    // 100MHz clock: T = 10ns
    always #5 clk = ~clk;

    always @(posedge clk or negedge rst_n) begin
        if (!rst_n)
            cycle_count <= 0;
        else
            cycle_count <= cycle_count + 1;
    end

    initial begin
        clk   = 0;
        rst_n = 0;
        cin_a = 8'd0;
        cin_b = 8'd0;
        c_in  = 1'b0;
        stall = 1'b0;
        flush = 1'b0;

        #20 rst_n = 1;
        #10;

        // Normal operation: 0x12 + 0x34 = 0x46
        cin_a = 8'h12;
        cin_b = 8'h34;

        // Cycle 10: stall for 2 cycles (stage 2)
        wait(cycle_count >= 10);
        stall = 1'b1;
        repeat(2) @(posedge clk);
        stall = 1'b0;

        // Cycle 15: flush (stage 3)
        wait(cycle_count >= 15);
        flush = 1'b1;
        @(posedge clk);
        flush = 1'b0;

        #100;
        $finish;
    end
endmodule
\end{lstlisting}
